\chapter{Reprise après sinistre et exigences réglementaires}\label{ch:nw:disaster-recovery-and-regulatory-requirements}

Lorsque Superstorm Sandy atteignit New York en octobre 2012, les bourses nationales américaines de valeurs mobilières «~ont décidé conjointement de ne pas ouvrir
pour la négociation les 29 et 30 octobre~». Le test annuel de la négociation depuis les sites de secours de l'industrie des valeurs mobilières avait eu lieu le
27 octobre, deux jours plus tôt, et n'avait rien relevé qui aurait empêché les marchés d'ouvrir sur leurs systèmes de secours~; beaucoup
firmes, la SEC l'a écrit plus tard, «~restaient mal à l'aise à l'idée de basculer~». Deux ans plus tard, la Regulation SCI exigeait que les
plans de reprise des bourses soient raisonnablement conçus pour rétablir leurs systèmes critiques dans les deux heures qui suivent
une perturbation de grande ampleur, et exigeait que les membres désignés par chaque plateforme testent leur site de secours avec elle au moins une fois par an.

Le livre~13 présente la bascule (failover), la réplication primaire--secours et le split brain~; le livre~6, le risque opérationnel~; le livre~11, le kill
switch~; le livre~15, la réponse aux incidents. Ce chapitre traite des sites et des règles~: où une firme place son site de reprise, comment les bourses
testent le leur avec leurs membres, sur quelles règles de résilience une société de trading est évaluée, et comment une conception de reprise se traduit en
un délai de reprise et un point de reprise. La simulation de reprise est signalée comme telle~; les règles sont citées textuellement.

\section{Sites secondaires et leur implantation}

\begin{definition}[Centre de données principal, site de reprise après sinistre]\label{def:nw:disaster-recovery-and-regulatory-requirements:sites}
Le \emph{centre de données principal}\index{centre de données principal} d'une firme est le site depuis lequel elle négocie normalement~; son
\emph{site de reprise après sinistre}\index{site de reprise après sinistre} est un second site, assez éloigné pour ne pas partager les
sinistres probables du site principal, et depuis lequel la firme peut reprendre la négociation lorsque ce dernier est perdu.
\end{definition}

Deux distances s'opposent. Le site de reprise doit être hors des aléas susceptibles d'emporter le site principal, une tempête, une
inondation, une panne régionale d'électricité~; et il doit être assez proche pour que sa réplication soit peu coûteuse et que ses connexions aux plateformes
soient utilisables. Les choix des bourses montrent l'ordre de grandeur~: les marchés américains de BrokerTec chez CME tournent à Secaucus et se reprennent à Aurora, ses
marchés européens tournent à Slough et se reprennent à Secaucus~; l'appariement au comptant de LSEG tourne à Slough et se reprend à Secaucus
(chapitres~23 à~25). Une reprise à travers un continent ou un océan, et non à travers le fleuve.

\begin{table}[ht]
\centering\small
\begin{tabular}{@{}lrl@{}}
\toprule
Site de reprise candidat & Depuis Secaucus NY5 (km) & Hors d'un rayon d'aléa de 100 km \\
\midrule
Secaucus NY4 (même campus) & 0.3 & non \\
Carteret & 26.0 & non \\
Mahwah & 34.2 & non \\
Aurora & 1\,191.1 & oui \\
\bottomrule
\end{tabular}
\caption{Sites de reprise candidats pour une société de trading négociant depuis Secaucus NY5~: \omterm{def:nw:the-north-american-map:geodesic}{distance géodésique} (coordonnées du chapitre~10) face à un
rayon d'aléa de \qty{100}{\kilo\metre} (une hypothèse). Données~: \texttt{nw\_dr.separations()}.}\label{tab:nw:disaster-recovery-and-regulatory-requirements:sites}
\end{table}

Les sites du New Jersey sont tous compris dans toute tempête qui fermerait Secaucus. Aurora ne l'est pas, et coûte ce que coûte la distance~: la réplication synchrone
vers ce site ajouterait \qty{17.4}{\milli\second} d'aller-retour à chaque écriture, pour un \omterm{def:nw:the-north-american-map:factor}{facteur de tracé} de 1.5, contre
\qty{0.5}{\milli\second} vers Mahwah~; la firme réplique donc de façon asynchrone et accepte un point de reprise.

\section{Tests de bascule des bourses}

\begin{definition}[Plan de continuité d'activité, test de bascule de la bourse]\label{def:nw:disaster-recovery-and-regulatory-requirements:test}
Un \emph{plan de continuité d'activité}\index{plan de continuité d'activité} est l'ensemble documenté de dispositions, de personnes et
procédures permettant de poursuivre ou de reprendre ses activités critiques après une perturbation. Un
\emph{test de bascule de la bourse}\index{test de bascule de la bourse} est un exercice programmé au cours duquel une bourse fait tourner ses marchés depuis son site de reprise ou ses composants de secours, et ses
membres se connectent, négocient et vérifient depuis les leurs.
\end{definition}

Les tests existent parce que Sandy a montré qu'un secours testé n'est pas forcément un secours fiable. La Regulation SCI les rend obligatoires pour les membres
désignés par chaque bourse, au moins une fois tous les douze mois, de façon coordonnée à l'échelle de l'industrie. L'industrie des contrats à terme organise la sienne~:
Le test 2025 de la FIA, le samedi 25 octobre avec un pré-test en septembre, a vu plus d'une centaine d'organisations vérifier
la connectivité aller-retour et la reprise des traitements entre les sites principal et de secours (l'\cref{dat:nw:disaster-recovery-and-regulatory-requirements:tests}).
Les bourses basculent aussi régulièrement~: les \omterm{def:nw:futures-connectivity:msgw}{passerelles de segment de marché} de CME passent sur leurs composants de secours puis reviennent, lors de fenêtres hebdomadaires.

\begin{dated}{2026-09}{Tests de bascule et sites de reprise des bourses}\label{dat:nw:disaster-recovery-and-regulatory-requirements:tests}
Test de reprise après sinistre de la FIA~: 25 octobre 2025 (pré-test le 27 septembre 2025), plus de 100 participants. \omterm{def:nw:futures-connectivity:msgw}{Passerelles de segment de marché} de CME~:
fenêtres de bascule hebdomadaires vers les composants de secours. BrokerTec U.S.~: NY5, reprise à Aurora, \omterm{def:nw:disaster-recovery-and-regulatory-requirements:objectives}{objectif de délai de reprise} de 4~heures~; BrokerTec
EU~: LD4.2, reprise à NY5, 2~heures. Appariement au comptant de LSEG~: LD4, reprise à NY6.
\end{dated}

\section{Règles de résilience sur lesquelles les firmes sont évaluées}

\begin{definition}[Résilience opérationnelle, service métier important, tolérance d'impact]\label{def:nw:disaster-recovery-and-regulatory-requirements:resilience}
La \emph{résilience opérationnelle}\index{résilience opérationnelle} est la capacité d'une firme à prévenir les perturbations, à s'y adapter, à y répondre, à s'en rétablir et
à en tirer les enseignements des perturbations opérationnelles. Un \emph{service métier important}\index{service métier important} est un service que la firme
fournit à ses clients ou au marché, et dont la perturbation pourrait leur causer un préjudice intolérable ou menacer la firme ou l'intégrité du marché~; sa
\emph{tolérance d'impact}\index{tolérance d'impact} est la perturbation maximale tolérable pour ce service, fixée par la firme, généralement exprimée en durée.
\end{definition}

\begin{dated}{2026-09}{Règles de résilience, d'après les textes des régulateurs}\label{dat:nw:disaster-recovery-and-regulatory-requirements:rules}
\begin{itemize}
\item Regulation SCI (bourses américaines et autres entités SCI)~: des capacités de secours «~suffisamment résilientes et géographiquement diversifiées~»
  et «~raisonnablement conçues pour permettre la reprise de la négociation dès le jour ouvré suivant et la reprise en deux heures des systèmes SCI critiques
  à la suite d'une perturbation de grande ampleur~»~; les membres désignés testent avec l'entité au moins une fois tous les 12 mois.
\item RTS~6, article~14 (firmes de trading algorithmique dans l'UE)~: dispositifs documentés de continuité d'activité par plateforme, scénarios
  y compris la perte de centres de données et de fournisseurs, le transfert vers un site de secours le cas échéant, l'arrêt ordonné~; revus et
  testés chaque année.
\item FCA PS21/3 (Royaume-Uni)~: \omterm{def:nw:disaster-recovery-and-regulatory-requirements:resilience}{services métier importants} et \omterm{def:nw:disaster-recovery-and-regulatory-requirements:resilience}{tolérances d'impact} fixés au 31 mars 2022~; capacité à rester dans ces limites d'ici le
  31 mars 2025. DORA (UE, Regulation 2022/2554)~: s'applique à partir du 17 janvier 2025.
\end{itemize}
\end{dated}

Les règles diffèrent selon ceux qu'elles engagent, mais elles partagent une même forme. La bourse doit rétablir ses systèmes critiques dans les deux heures~; la société de trading
doit disposer de dispositifs lui permettant de reprendre ou d'arrêter ses activités dans l'ordre, testés chaque année~; et au Royaume-Uni, la firme elle-même désigne les
services qui comptent et la perturbation qu'elle acceptera de tolérer, puis montre qu'elle peut rester dans cette limite. La conception de reprise d'une société de trading est
donc jugée sur deux chiffres qu'elle se fixe elle-même, un délai de reprise et un point de reprise, au regard de la tolérance qu'elle a déclarée.

\section{Objectifs de reprise et runbooks}

\begin{definition}[Objectif de délai de reprise (RTO), objectif de point de reprise (RPO)]\label{def:nw:disaster-recovery-and-regulatory-requirements:objectives}
Un \emph{objectif de délai de reprise}\index{objectif de délai de reprise} (RTO) est la durée maximale qu'un système peut mettre, à compter de sa perte, pour
redevenir opérationnel~; un \emph{objectif de point de reprise}\index{objectif de point de reprise} (RPO) est la donnée la plus récente que la firme peut perdre, mesurée comme
le temps écoulé entre la dernière modification détenue par le site de reprise et l'instant de la perte.
\end{definition}

Une reprise est un runbook~: détecter la perte et décider de la bascule, transférer les sessions vers les adresses de reprise, restaurer ou démarrer ce qui
n'est pas en marche, charger les positions et les données de référence, rétablir et vérifier les sessions, réconcilier les ordres et les positions avec les
drop copies des bourses, reprendre. Chaque étape a une durée qui varie~; le délai de reprise en est la somme.

\omcode{code/firm/drplan/firm_drplan.py}{61}{84}{Délai de reprise comme somme des étapes d'un runbook, point de reprise selon le mode de réplication et coût de la réplication synchrone.}\label{lst:nw:disaster-recovery-and-regulatory-requirements:rto}

\begin{proposition}[Ce qui fixe chaque objectif]\label{prop:nw:disaster-recovery-and-regulatory-requirements:objectives}
Le point de reprise ne dépend que de la réplication~: nul si la réplication est synchrone, le retard de réplication plus le délai
unidirectionnel si elle est asynchrone, jusqu'à l'intervalle entre deux instantanés si la reprise se fait à partir d'instantanés. Le délai
de reprise ne dépend que du runbook~: il est au moins égal à la somme des durées minimales des étapes, quelle que soit la réplication.
\end{proposition}

\begin{proof}
Les données que le site de reprise ne détient pas au moment de la perte du principal sont perdues~; leur volume est fixé par le moment où elles ont été envoyées, c'est-à-dire par la réplication.
Les étapes du runbook s'enchaînent après la perte, et aucune ne peut durer moins que son minimum~; aucun des deux objectifs ne peut donc
être obtenu avec l'outil de l'autre.
\end{proof}

La firme du chapitre, qui opère depuis Secaucus avec un site de reprise à Aurora, compare deux conceptions
(\cref{fig:nw:disaster-recovery-and-regulatory-requirements:rto}). La conception à chaud maintient des serveurs en fonctionnement à Aurora avec une
réplication asynchrone (un retard de \qty{50}{\milli\second}) et des sessions prêtes~: son runbook comporte cinq étapes, de médianes 15, 10, 20, 10 et 5 minutes.
La conception à tiède maintient le site et quelques serveurs, restaure à partir d'instantanés pris toutes les quinze minutes, et compte six étapes, dont 45
minutes pour démarrer et restaurer et 30 pour charger les données. Toutes les durées sont des hypothèses, lognormales autour de leurs médianes.

\omcode{code/networks/28-disaster-recovery-and-regulatory-requirements/python/nw_dr.py}{24}{34}{Les deux conceptions~: runbooks, réplication et coût annuel.}\label{lst:nw:disaster-recovery-and-regulatory-requirements:plans}

\begin{omfigure}
\begin{tikzpicture}
\begin{axis}[width=11.5cm,height=5.8cm,xmin=0,xmax=300,ymin=0,ymax=0.36,
  xlabel={délai de reprise (minutes)},ylabel={part par classe de 10 min},
  tick label style={font=\footnotesize},label style={font=\small},grid=major,grid style={gray!25},
  legend pos=north east,legend style={font=\footnotesize},table/col sep=comma]
\addplot[omDef,thick,fill=omDef!25,const plot mark mid,area legend] table[x=minutes,y=hot]{figdata/networks/28-disaster-recovery-and-regulatory-requirements/rto.csv} \closedcycle;
\addplot[omThm,thick,fill=omThm!20,fill opacity=0.6,const plot mark mid,area legend] table[x=minutes,y=warm]{figdata/networks/28-disaster-recovery-and-regulatory-requirements/rto.csv} \closedcycle;
\draw[omMeth,very thick,dashed] (axis cs:120,0) -- (axis cs:120,0.36);
\node[text=omMeth,font=\footnotesize,anchor=west] at (axis cs:122,0.33) {deux heures};
\legend{site à chaud,site à tiède}
\end{axis}
\end{tikzpicture}

\omcaption{Simulation~: les distributions des délais de reprise des conceptions à chaud et à tiède, face à la ligne des deux heures. La conception à chaud reprend
en deux heures dans 99.2~\% des reprises simulées (médiane 64 minutes)~; la conception à tiède dans 20.3~\% (médiane 148). Données~:
\texttt{fig\_dr.py}, \texttt{nw\_dr.rto\_hist()}.}\label{fig:nw:disaster-recovery-and-regulatory-requirements:rto}
\end{omfigure}

La conception à chaud reprend en deux heures dans 99.2~\% des reprises simulées, avec une médiane de 64 minutes et un 95\up{e}~centile à 96~; son
point de reprise est de \qty{56}{\milli\second}, soit le retard plus le délai unidirectionnel jusqu'à Aurora. La conception à tiède respecte les deux heures dans 20.3~\%, avec une
médiane de 148 minutes et un 95\up{e}~centile à 237, et peut perdre quinze minutes de données. Si le site à chaud coûte 70~\% du
coût annuel de l'infrastructure du principal et le site à tiède 25~\% (la pile de l'exemple du chapitre~27, 3.19 millions de dollars par an~; hypothèses), l'écart atteint
1.44 million de dollars par an~: le prix de la ligne des deux heures pour cette firme.

\begin{method}[Concevoir une reprise pour une firme de trading]\label{met:nw:disaster-recovery-and-regulatory-requirements:design}
\begin{enumerate}
\item Nommer les \omterm{def:nw:disaster-recovery-and-regulatory-requirements:resilience}{services métier importants} et fixer une \omterm{def:nw:disaster-recovery-and-regulatory-requirements:resilience}{tolérance d'impact} pour chacun~; en déduire les \omterm{def:nw:disaster-recovery-and-regulatory-requirements:objectives}{objectifs de délai de reprise} et de point de reprise.
\item Choisir un site de reprise à l'écart des risques du principal et proche des sites de reprise propres aux plateformes~; calculer ce que coûterait
  une réplication synchrone en termes de latence, et répliquer de manière asynchrone si ce coût est trop élevé.
\item Écrire le runbook sous forme d'étapes chronométrées~; mesurer chaque étape lors des tests, et non dans des estimations.
\item Simuler le délai de reprise par rapport à l'objectif~; réduire les étapes qui dominent, ou changer de conception.
\item Participer aux tests de bascule des bourses et de la profession, et en mener un propre au moins une fois par an~; consigner les résultats
  comme preuves destinées au régulateur.
\end{enumerate}
\end{method}

\section{Tutoriel~: la règle des deux heures}\label{tut:nw:disaster-recovery-and-regulatory-requirements:tut}

\begin{tutorial}
\textbf{Objectif.} Placer un site de reprise, simuler les conceptions à chaud et à tiède, et les comparer à la ligne des deux heures et à leur coût.
\textbf{État final~:} \cref{fig:nw:disaster-recovery-and-regulatory-requirements:rto} et
\cref{tab:nw:disaster-recovery-and-regulatory-requirements:sites}.
\begin{tutsteps}
\item \textbf{Sites.} \texttt{firm\_drplan.separation\_km} et \texttt{outside\_hazard} sur les coordonnées du chapitre~10.
\item \textbf{Plans.} \texttt{Plan} et \texttt{Step} (\cref{lst:nw:disaster-recovery-and-regulatory-requirements:plans})~;
  \texttt{rto\_samples}, \texttt{p\_within} et \texttt{rpo} (\cref{lst:nw:disaster-recovery-and-regulatory-requirements:rto}).
\item \textbf{Calendrier.} \texttt{load\_calendar} et \texttt{tested\_within} pour le test annuel.
\end{tutsteps}
\textbf{Pour aller plus loin.} Faire fonctionner certaines étapes en parallèle (la réconciliation pendant que les sessions se rétablissent) et voir laquelle
des deux conceptions voit sa médiane bouger~; ajouter la probabilité que la décision de basculer soit retardée, comme elle l'a été avant Sandy.
\end{tutorial}

\section{Construction~: le plan de reprise après sinistre}\label{bld:nw:disaster-recovery-and-regulatory-requirements:drplan}

\begin{build}
\bfield{Rôle} Un plan de reprise sous forme de données, avec son délai et son point de reprise, des contrôles par rapport aux cibles et un calendrier de tests~: les
lignes de résilience du plan du chapitre~29.
\bfield{Interface} \texttt{firm\_drplan}~: \texttt{Step}, \texttt{Plan}, \texttt{separation\_km}, \texttt{outside\_hazard},
\texttt{rto\_samples}, \texttt{p\_within}, \texttt{rpo}, \texttt{sync\_penalty\_ms}, \texttt{tested\_within}, \texttt{load\_calendar}.
\bfield{Règles} Les sites proviennent de la table des sites avec leurs sources~; les dates des tests sont des lignes datées~; les durées des runbooks sont des hypothèses énoncées,
remplacées par des durées mesurées après chaque test.
\bfield{Tests d'acceptation} \texttt{code/firm/drplan/tests/}~: un runbook déterministe, le point de reprise selon le mode, les séparations, et le
contrôle du test annuel.
\bfield{Pour aller plus loin} Étapes en parallèle~; défaillances corrélées du principal et d'un fournisseur~; le délai de reprise propre à la bourse ajouté à celui de la firme.
\end{build}

\begin{omsources}
\item 17 CFR 242.1001 et 242.1004 (Regulation SCI)~; SEC Releases 34-69077 (2013) et 34-73639 (2014).
\item Commission Delegated Regulation (EU) 2017/589, Article 14~; FCA PS21/3~; ESMA on DORA.
\item FIA, 2025 disaster recovery test.
\end{omsources}

\section{Exercices}

\begin{exercise}[$\star$]\label{exo:nw:disaster-recovery-and-regulatory-requirements:1}
Quelles capacités de sauvegarde une entité SCI doit-elle être conçue pour garantir après une perturbation à grande échelle~?
\end{exercise}

\begin{exercise}[$\star$]\label{exo:nw:disaster-recovery-and-regulatory-requirements:2}
Quelle est la différence entre un \omterm{def:nw:disaster-recovery-and-regulatory-requirements:objectives}{objectif de délai de reprise} et un \omterm{def:nw:disaster-recovery-and-regulatory-requirements:objectives}{objectif de point de reprise}~?
\end{exercise}

\begin{exercise}[$\star$]\label{exo:nw:disaster-recovery-and-regulatory-requirements:3}
Pourquoi Mahwah est-il un mauvais site de reprise pour une firme qui opère depuis Secaucus~?
\end{exercise}

\begin{exercise}[$\star\star$]\label{exo:nw:disaster-recovery-and-regulatory-requirements:4}
En utilisant la \cref{prop:nw:disaster-recovery-and-regulatory-requirements:objectives}, pourquoi une réplication plus rapide ne raccourcit-elle pas le
délai de reprise de la conception à tiède~?
\end{exercise}

\begin{exercise}[$\star\star$]\label{exo:nw:disaster-recovery-and-regulatory-requirements:5}
Que doit exiger la RTS 6 des dispositifs de continuité d'activité d'une société de trading concernant les ordres en suspens~?
\end{exercise}

\begin{exercise}[$\star\star$]\label{exo:nw:disaster-recovery-and-regulatory-requirements:6}
L'entreprise a effectué son dernier test de bascule le 25 octobre 2025. Ce test est-il toujours valable dans les douze mois le 28 septembre 2026, et le 1\up{er} novembre 2026~?
\end{exercise}

\begin{exercise}[$\star\star\star$]\label{exo:nw:disaster-recovery-and-regulatory-requirements:7}
\emph{Programmation.} Avec \texttt{firm\_drplan}, de combien la réduction de moitié de l'étape de restauration de la conception à tiède (de 45 à 22.5 minutes) augmenterait-elle
sa probabilité de reprendre dans les deux heures~?
\end{exercise}

\begin{exercise}[$\star\star\star$]\label{exo:nw:disaster-recovery-and-regulatory-requirements:8}
\emph{Trouver la faille.} «~Nous avons testé notre site de secours le mois dernier, et il a fonctionné~; nous basculerons donc dessus si une tempête survient.~»
\end{exercise}

\section{Problème~: La règle des deux heures}

\begin{problem}[{Problème du week-end --- reprise à chaud contre reprise à tiède}]\label{pb:nw:disaster-recovery-and-regulatory-requirements:1}
Une firme négocie depuis Secaucus NY5 et doit choisir une conception de reprise. Elle a fait sienne la ligne des deux heures
des bourses comme \omterm{def:nw:disaster-recovery-and-regulatory-requirements:objectives}{objectif de délai de reprise}. Utiliser le modèle du chapitre.

\medskip\noindent\textbf{Partie I --- Le site.}
\begin{enumerate}
\item À quelle distance de NY5 se trouve chaque site candidat, et lesquels sont situés hors d'un rayon de risque de 100~kilomètres~?
\item Qu'ajouteraient la réplication synchrone vers Aurora et vers Mahwah à chaque écriture~?
\item Pourquoi les bourses se rétablissent-elles si loin~?
\item Quel point de reprise la réplication asynchrone vers Aurora donne-t-elle~?
\end{enumerate}

\medskip\noindent\textbf{Partie II --- Les conceptions.}
\begin{enumerate}[resume]
\item Quelle est la probabilité de reprise à chaud dans les deux heures, sa médiane et son 95\up{e}~centile~?
\item Et celle de la conception à tiède~?
\item Quelle quantité de données la conception à tiède peut-elle perdre~?
\item Quelles étapes dominent le temps de la conception à tiède~?
\end{enumerate}

\medskip\noindent\textbf{Partie III --- Les règles.}
\begin{enumerate}[resume]
\item Qu'impose la Regulation SCI aux bourses et à leurs membres désignés~?
\item Qu'impose la RTS 6 aux sociétés de trading~?
\item Qu'imposaient les règles britanniques au 31 mars 2025~?
\item Qu'a montré Sandy que les tests seuls ne montraient pas~?
\end{enumerate}

\medskip\noindent\textbf{Partie IV --- Le verdict.}
\begin{enumerate}[resume]
\item Énoncer le \emph{résultat nommé}~: la probabilité que les conceptions à chaud et à tiède reprennent dans les deux heures, et le coût annuel de la
  différence.
\item Quelles données sont publiées, et lesquelles sont supposées~?
\item La conception à chaud vaut-elle 1.44~million de dollars par an~?
\item Qu'est-ce qui permettrait à la conception à tiède de respecter la ligne plus souvent~?
\item Quelles preuves un régulateur demanderait-il~?
\item Comment la firme devrait-elle participer aux tests des bourses~?
\item Où cela s'insère-t-il dans le plan du chapitre~29~?
\item En une phrase~: que procure une conception de reprise~?
\end{enumerate}
\end{problem}

\section{Questions d'entretien}

\begin{interviewq}[$\star$ \iqroles{développeur}]\label{iq:nw:disaster-recovery-and-regulatory-requirements:1}
Que sont le RTO et le RPO, et qu'est-ce qui détermine chacun~?
\end{interviewq}

\begin{interviewq}[$\star\star$ \iqroles{développeur}]\label{iq:nw:disaster-recovery-and-regulatory-requirements:2}
Où placeriez-vous le site de reprise d'une firme qui négocie à Secaucus, et pourquoi~?
\end{interviewq}

\begin{interviewq}[$\star\star$ \iqroles{développeur}]\label{iq:nw:disaster-recovery-and-regulatory-requirements:3}
Comment évitez-vous le split brain lors de la bascule d'un système de trading~?
\end{interviewq}

\begin{interviewq}[$\star\star$ \iqroles{trader}]\label{iq:nw:disaster-recovery-and-regulatory-requirements:4}
Le site principal est perdu alors que des positions sont ouvertes. Que se passe-t-il dans les trente premières minutes~?
\end{interviewq}

\begin{interviewq}[$\star\star$ \iqroles{développeur, chercheur}]\label{iq:nw:disaster-recovery-and-regulatory-requirements:5}
Comment décideriez-vous entre un site de reprise à chaud et un site de reprise à tiède~?
\end{interviewq}

\begin{interviewq}[$\star\star\star$ \iqroles{développeur, chercheur}]\label{iq:nw:disaster-recovery-and-regulatory-requirements:6}
Concevez les dispositifs de continuité d'activité d'une société de trading algorithmique active aux États-Unis et dans l'UE, de façon
qu'ils satisfassent aux règles auxquelles elle sera soumise.
\end{interviewq}
